Bài tập lớn môn mạng nơ-ron sâu
Môn học: Deep Neural Networks
Hình thức: nhóm 2–3 sinh viên
Môi trường: Google Colab (GPU T4) · PyTorch
1. Bối cảnh
Một công ty muốn xây dựng bộ phân loại ảnh cho 10 lớp đồ vật và động vật (CIFAR-10),
nhưng chỉ có 5.000 ảnh đã gán nhãn để huấn luyện. Huấn luyện từ đầu với lượng dữ liệu
này cho kết quả kém. Nhiệm vụ của nhóm là:
1. Hiểu và giải thích được vì sao từng thành phần khi huấn luyện (optimizer, learning rate,
khởi tạo, regularization) ảnh hưởng tới kết quả.
2. Dùng transfer learning và fine-tuning để vượt qua mô hình baseline do giảng viên cung
cấp.
3. Thể hiện rằng nhóm biết debug, theo dõi và gỡ lỗi quá trình huấn luyện.
4. Đóng gói mô hình tốt nhất thành một demo có giao diện.
2. Dữ liệu
Tập Nguồn Số ảnh
train lấy từ CIFAR-10 train, 500 ảnh/lớp 5.000
val lấy từ CIFAR-10 train, 500
ảnh/lớp, không trùng train
5.000
test CIFAR-10 test chính thức 10.000
Chỉ số ảnh của train và val được sinh cố định bởi hàm make_splits(seed=42)
trong notebooks/baseline.ipynb. Mọi nhóm phải dùng đúng hàm này.
Luật bắt buộc:
● Không dùng thêm ảnh nào từ 45.000 ảnh CIFAR-10 train còn lại, kể cả ảnh không nhãn.
● Không dùng mô hình đã được pretrain hoặc fine-tune trên CIFAR-10/CIFAR-100.
● Cố định seed và báo cáo kết quả trung bình ± độ lệch chuẩn trên ít nhất 3 seed cho mô
hình cuối cùng.
3. Baseline
notebooks/baseline.ipynb
4. Nội dung cần thực hiện
Phần A–D là bắt buộc. Phần E chọn ít nhất 1 hướng.
Phần A: Thuật toán tối ưu (SGD, Momentum, RMSProp, Adam)
A1. Tự cài đặt 4 optimizer dưới dạng lớp kế thừa torch.optim.Optimizer, không gọi
torch.optim.SGD/Adam/RMSprop bên trong. Viết công thức cập nhật trong báo cáo,
kèm ý nghĩa từng thành phần (vận tốc, trung bình bình phương gradient, bias correction, ε).
A2. Kiểm chứng bản tự cài: chạy song song với optimizer của PyTorch trên cùng mô hình
nhỏ, cùng khởi tạo, cùng dữ liệu trong 100 bước. Sai khác tham số phải < 1e-5.
Phần B: Learning rate, khởi tạo trọng số, regularization
B1. LR range test: tăng LR theo cấp số nhân từ 1e-7 đến 10 trong một epoch, vẽ loss theo
LR và giải thích cách chọn LR từ đồ thị.
B2. LR schedule: so sánh ít nhất 3 schedule (constant, step decay, cosine, warmup + cosine,
OneCycle, …). Vẽ LR thực tế theo từng bước để chứng minh schedule chạy đúng như mong
muốn.
B3. Khởi tạo trọng số: trên một CNN nhỏ huấn luyện từ đầu (≥ 6 lớp conv, không
BatchNorm), so sánh: khởi tạo toàn 0, N(0, 1), Xavier/Glorot, He/Kaiming. Với mỗi cách,
vẽ histogram activation và độ lớn gradient của từng lớp ở bước đầu tiên. Giải thích kết quả
bằng lý thuyết phương sai.
B4. Regularization: ablation trên mô hình chính, gồm có/không Dropout, có/không
BatchNorm (với CNN tự xây), weight decay, data augmentation. Lập bảng train acc, val acc
và khoảng cách train–val cho từng cấu hình.
Phần C: Transfer learning và fine-tuning
C1. Sử dụng bốn phương pháp đã học trên cùng dữ liệu huấn luyện:
1. Huấn luyện từ đầu (scratch), ResNet-18 khởi tạo ngẫu nhiên.
2. Linear probe (đóng băng backbone).
3. Full fine-tuning với learning rate nhỏ.
4. Fine-tuning với discriminative learning rate (dùng parameter group: LR nhỏ cho lớp
đầu, LR lớn cho head).

C2. Đường cong theo lượng dữ liệu: chạy cả 4 chiến lược với 1%, 10% và 100% tập
train (50 / 500 / 5.000 ảnh). Vẽ accuracy theo lượng dữ liệu (trục x dạng log). Phân tích: ở
mức dữ liệu nào thì scratch bắt đầu đuổi kịp, và vì sao.
C3. BatchNorm khi fine-tune: so sánh việc để BN ở chế độ train và đóng băng running
statistics, với batch size lớn (128) và nhỏ (8). Giải thích kết quả.
C5. Trả lời các câu hỏi sau:
● Khi nào nên chọn linear probe thay vì full fine-tuning?
● Vì sao một learning rate tốt khi train từ đầu lại phá hỏng mô hình pretrained?
● Nhóm quyết định số lớp cần unfreeze như thế nào?
Phần D: Huấn luyện và gỡ lỗi
D1. Sanity check bắt buộc trước khi huấn luyện thật, trình bày bằng số liệu và hình trong
báo cáo:
● Loss ban đầu ≈ ln(10) ≈ 2,30. Nếu không, giải thích vì sao.
● Overfit được một batch duy nhất (loss → gần 0).
● Hiển thị một batch sau augmentation kèm nhãn.
● Shuffled-label test: huấn luyện với nhãn bị xáo trộn, val accuracy phải ≈ 10%.
D2. Nhật ký gỡ lỗi: ghi lại ít nhất 2 lỗi thật nhóm gặp phải trong quá trình làm bài, gồm triệu
chứng, cách phát hiện và cách sửa.
Phần E: Nghiên cứu mở rộng (chọn ít nhất 1)
Phần này đánh giá khả năng nghiên cứu vượt ra ngoài nội dung môn học. Mỗi hướng cần có:
tóm tắt bài báo hoặc phương pháp gốc, thí nghiệm có đối chứng và phân tích.
● PEFT: LoRA hoặc Adapter cho một backbone lớn hơn (ViT-S, ConvNeXt-T) bằng thư
viện peft. So sánh accuracy và số tham số trainable với full fine-tuning.
● Optimizer hiện đại: AdamW (so với Adam + L2), Lion, SAM, Lookahead, … Giải thích
cơ chế và kiểm chứng.
● Regularization nâng cao: Mixup, CutMix, label smoothing, RandAugment, stochastic
depth.
● Knowledge distillation: teacher lớn → student ResNet-18.
● Domain shift / negative transfer: đem pipeline sang một tập khác lệch miền ImageNet
(EuroSAT, ảnh y tế MedMNIST, …) và phân tích.
● Độ tin cậy và giải thích mô hình: calibration (ECE, temperature scaling), Grad-CAM,
phân tích lỗi theo lớp.
● Hướng khác do nhóm đề xuất, cần được giảng viên duyệt trước tuần 3.
5. Sản phẩm nộp
1. Báo cáo (PDF, 15–25 trang, không tính phụ lục) theo report_template.md.
2. Mã nguồn: notebook(s) chạy được trên Colab từ đầu đến cuối, kèm file cấu hình và log thí
nghiệm. Ghi rõ lệnh hoặc cell để tái lập kết quả cuối.
3. Checkpoint mô hình tốt nhất (link Google Drive hoặc Hugging Face).
4. Demo: giao diện bất kỳ (Gradio, Streamlit, web app, …) cho phép tải ảnh lên và nhận topk dự đoán kèm xác suất. Có mẫu khởi đầu tại demo/app_gradio.py. Khuyến khích
deploy lên Hugging Face Spaces.
5. Thuyết trình 10 phút + 5 phút vấn đáp, có demo trực tiếp.